\section*{Capítulo \ref{ch:b3:lab-techniques-3} --- Técnicas de laboratório III: a prática da pesquisa}

\begin{solution}{exo:b3:lab-techniques-3:1}
$(0.50/1013)^3 = \num{1.2e-10}$. Os vazamentos, o gás adsorvido no vidro e a liberação de gases pela graxa e pelos septos
deixam muito mais que isso; vidraria seca e quente e boas juntas importam mais que
ciclos extras.
\end{solution}

\begin{solution}{exo:b3:lab-techniques-3:2}
$\qty{20.0}{mmol}/\qty{1.48}{mol/L} = \qty{13.5}{mL}$.
\end{solution}

\begin{solution}{exo:b3:lab-techniques-3:3}
$120 + 10 \times 1.00782503 + 55.93493633 - 0.00054858 = \qty{186.01264}{u}$. Erro: $10^6 \times (186.0129 - 186.01264)/186.01264 = +1.4$ ppm,
bem dentro de 5 ppm.
\end{solution}

\begin{solution}{exo:b3:lab-techniques-3:4}
Primária: o artigo de revista, a patente, a tese. Secundária: o artigo de revisão, o
manual de constantes. Terciária: o livro didático.
\end{solution}

\begin{solution}{exo:b3:lab-techniques-3:5}
$M = \qty{194.0}{g/mol}$: C $96.0/194.0 = 49.48$\,\%, H $10.0/194.0 = 5.15$\,\%, N $56.0/194.0 = 28.87$\,\%. Diferenças de 0.17, 0.07 e 0.16
ponto: dentro de 0.4, a análise apoia a fórmula.
\end{solution}

\begin{solution}{exo:b3:lab-techniques-3:6}
$\qty{170.0}{mg}/\qty{212.0}{g/mol} = \qty{0.802}{mmol}$; $c = \qty{0.802}{mmol}/\qty{0.54}{mL} = \qty{1.5}{mol/L}$ (dois algarismos significativos, fixados pela
leitura do volume).
\end{solution}

\begin{solution}{exo:b3:lab-techniques-3:7}
Médias de 73.0 e 77.0\,\%; os dois desvios-padrão valem $\sqrt{10/4} = 1.58$; $s_p = 1.58$. $t = 4.0/(1.58\sqrt{2/5}) = 4.0$, maior que
2.31: a melhoria é significativa ao nível de 95\,\%.
\end{solution}

\begin{solution}{exo:b3:lab-techniques-3:8}
$F = (0.80/0.30)^2 = 7.1 > 5.05$: o primeiro método é significativamente menos preciso.
\end{solution}

\begin{solution}{exo:b3:lab-techniques-3:9}
Compre frascos pequenos com inibidor, date cada um no recebimento e na abertura, armazene-os
fechados, no escuro e no frio; teste os peróxidos antes de qualquer destilação ou evaporação
e a intervalos fixos após a abertura; descarte os frascos vencidos, ou com teste
positivo, pela rota de resíduos perigosos, sem concentrá-los.
\end{solution}

\begin{solution}{exo:b3:lab-techniques-3:10}
Perigos: o hidrogênio liberado (incêndio, pressão), a reatividade do NaH com a água
e uma decomposição exotérmica conhecida da DMF com o hidreto de sódio, que pode sair de controle
sob aquecimento; a DMF é também tóxica para a reprodução. Medidas: substituir o solvente
(THF ou 2-metiltetra-hidrofurano) ou a base; se mantidos, começar em pequena escala, a baixa
temperatura, adicionando o NaH em porções e acompanhando a temperatura, em capela,
atrás de um anteparo, com saída por um borbulhador e um banho de resfriamento pronto; procedimento
escrito, uma segunda pessoa presente e, por último, o equipamento de proteção.
\end{solution}

\begin{solution}{exo:b3:lab-techniques-3:11}
Contribuições relativas: 0.50, 0.50, $0.010/20.0 = 0.05$, $0.010/15.0 = 0.067$ e 0.10\,\%; combinada,
$\sqrt{0.25 + 0.25 + 0.0025 + 0.0044 + 0.010} = 0.72$\,\%. As duas integrais dominam: uma melhor relação sinal/ruído e uma
correção cuidadosa da linha de base e da fase importam mais que a balança.
\end{solution}

\begin{solution}{exo:b3:lab-techniques-3:12}
$Q = (89.0 - 82.0)/(89.0 - 80.6) = 0.83 > 0.71$: o teste o sinaliza. Antes de rejeitá-lo, procure no caderno
uma causa (uma bureta mal lida, uma bolha de ar, outro lote de titulante); se encontrar uma,
rejeite-o e diga por quê; se não, mantenha-o, relate o teste ou repita a
medida.
\end{solution}

\begin{solution}{pb:b3:lab-techniques-3:1}
\textbf{1.} $\ce{FeCl2 + 2NaC5H5 -> Fe(C5H5)2 + 2NaCl}$.

\textbf{2.} O ciclopentadienil sódio é destruído pelo oxigênio e pela água, e o cloreto de
ferro(II) se oxida a ferro(III) ao ar úmido.

\textbf{3.} $\qty{5.00}{g}/\qty{126.8}{g/mol} = \qty{0.0394}{mol}$; $\times \qty{185.8}{g/mol} = \qty{7.33}{g}$; rendimento $5.86/7.33 = 80$\,\%.

\textbf{4.} $(0.10/1000)^3 = 10^{-12}$ (em teoria).

\textbf{5.} Aquecendo o dímero, que sofre uma reação de retro-Diels--Alder, e
destilando o monômero volátil; à temperatura ambiente, o monômero volta a se dimerizar
por uma reação de Diels--Alder, de modo que ele é mantido frio e usado imediatamente.

\textbf{6.} Cada porção libera hidrogênio e calor: adicioná-lo lentamente mantém sob controle a vazão
de gás e a temperatura.

\textbf{7.} \qty{186.0126}{u}; erro de $+1.4$ ppm.

\textbf{8.} C $120.0/185.8 = 64.58$\,\%, H $10.0/185.8 = 5.38$\,\%; encontrados 64.41 e 5.47: diferenças de $-0.17$ e $+0.09$, dentro de 0.4.

\textbf{9.} Os dez hidrogênios são equivalentes por simetria, e os anéis giram rapidamente.

\textbf{10.} Um único sinal, pois os dez carbonos são equivalentes.

\textbf{11.} Uma estrutura de raios X em monocristal.

\textbf{12.} Não: como complexo de 18 elétrons (\cref{ch:b3:organometallic-bonding}), ele é estável ao ar, ao contrário de seus
precursores.

\textbf{13.} Seus sinais não se sobrepõem aos do ferroceno, ele é um sólido puro, não volátil e
não higroscópico que pode ser pesado com exatidão, e é inerte em relação à
amostra.

\textbf{14.} A relaxação completa de todos os sinais entre as varreduras (um intervalo de vários $T_1$),
uma excitação uniforme, uma relação sinal/ruído suficiente e uma correção cuidadosa da fase e da
linha de base.

\textbf{15.} Ferroceno $\ce{C10H10Fe}$ \qty{185.8}{g/mol}; 1,3,5-trimetoxibenzeno $\ce{C9H12O3}$ \qty{168.0}{g/mol}.

\textbf{16.} $P_x = 3.942 \times (3/10) \times (185.8/168.0) \times (15.0/20.0) \times 0.999 = 0.980$: 98.0\,\%.

\textbf{17.} Relativa $\sqrt{0.50^2 + 0.50^2 + 0.05^2 + 0.067^2 + 0.10^2} = 0.72$\,\%; absoluta $0.72\,\% \times 98.0 = 0.7$ ponto percentual.

\textbf{18.} Sim, mas fracamente: uma impureza de 2\,\% de composição semelhante muda C e H menos
do que a análise consegue detectar; a qNMR é a medida de pureza mais informativa.

\textbf{19.} O hidreto de sódio (hidrogênio com a água, incêndio), o hidrogênio liberado, o THF e o
ciclopentadieno (altamente inflamáveis; o THF, \omterm{def:b3:lab-techniques-3:peroxide}{formador de peróxidos}), o craqueamento a quente do
dímero, o cloreto de ferro(II) (irritante) e o ferroceno (sólido inflamável, nocivo).

\textbf{20.} Hidreto de sódio: pesado como dispersão em óleo sob gás inerte, adicionado em
porções a uma solução agitada e resfriada, com o gás saindo por um borbulhador e nenhuma água
por perto. THF: seco e sem peróxidos, vindo de uma coluna, usado na capela sob nitrogênio.

\textbf{21.} Sob gás inerte e com resfriamento, por adição lenta de isopropanol, depois de
etanol, depois de água, esperando cada vez até que cesse o desprendimento de gás.

\textbf{22.} Uma onda reversível de um elétron, $\ce{Fe(C5H5)2+}/\ce{Fe(C5H5)2}$, com picos separados por cerca de 59 mV a \qty{25}{\celsius}
(\cref{ch:b3:electrode-kinetics}); o par é rápido, estável e quase insensível ao
solvente, o que faz dele uma referência interna conveniente.

\textbf{23.} O caderno: data, esquema, quantidades e origens, a \omterm{def:b3:lab-techniques-3:assessment}{avaliação de riscos}, o que
foi feito e observado, rendimentos e os nomes dos arquivos de \omterm{def:b3:lab-techniques-3:notebook}{dados brutos}. O \omterm{def:b3:lab-techniques-3:literature}{material
suplementar}: o procedimento completo e todos os dados de caracterização, com cópias dos
espectros.

\textbf{24.} Para que outros, e os próprios autores anos depois, possam reabri-los e reprocessá-los
qualquer que seja o programa de que disponham: localizáveis, acessíveis, interoperáveis e reutilizáveis.

\textbf{25.} A pureza por qNMR da amostra de ferroceno é de $98.0 \pm 0.7$\,\% (incerteza-padrão).
\end{solution}
